\section{Runtime}
\label{sec:runtime}

A machine of blocks needs what a single interface did not: time, input richer than a click or a drag, and a
loop that moves events between blocks while checking every step. This section specifies that loop and the input
model, and fixes the one design rule that keeps continuous motion from wandering.

\subsection{The tick}
\label{sec:tick}

The runtime is one loop. Each frame it collects inputs, turns them into actions through the bindings
(\cref{sec:bindings}), applies each action to its block, drains the event queue through the wiring with each
delivery stepping its target block, runs every stepped block's spec check, and renders. A tick is a pure
function of the machine state and the frame's inputs, so a session is a log of inputs and can be replayed
exactly, which is also how it is tested. Events emitted during a tick are processed within it, in a fixed order,
to the bound the machine declares; the clock is a block like any other, whose only action is \texttt{advance} and
whose event the machine wires to whatever it drives. This is the synchronous hypothesis of
Esterel and Lustre~\citep{berry1992esterel,halbwachs1991lustre}: a reaction is instantaneous and deterministic,
and causality is settled when the machine is checked, not when it runs. Harel's statecharts~\citep{harel1987statecharts}
and the Elm architecture~\citep{czaplicki2013elm} are the same discipline for interfaces, and ours differs only
in what steps a block.

\subsection{Bindings: the input model}
\label{sec:bindings}

Desktop software is mostly keys, modifiers and selection, not drags. None of that is new state; it is a mapping
from input events to actions, and it belongs in the spec rather than in code. A \texttt{bindings} section maps
a key or mouse event, with modifiers, to an action whose parameters come from the pointer, the focused item or
the selection (\cref{fig:bindings}). Modifiers change which action a gesture performs: Alt-drag performs
\texttt{copy} where drag performs \texttt{move} (Dragology's \texttt{alt-to-copy} and \texttt{drag-to-copy});
Shift-drag constrains a two-dimensional \texttt{vary} to one axis by fixing the other field; the wheel steps a
field; a double-click opens an item; Esc leaves a mode. The runtime owns the mapping and the blocks never see a
key, so the same machine runs under a different keyboard layout, or under a touch screen, by changing one
section.

\begin{figure}[t]
\begin{lstlisting}
bindings:
  Tab:           focus.next
  shift+Tab:     focus.prev
  ArrowRight:    board.move(card: focus.id, dx: 1, dy: 0)
  Delete:        board.remove(card: selection.items)
  ctrl+a:        selection.all
  Escape:        [selection.none, menu.close]
  click(card):   selection.only(card)
  ctrl+click:    selection.toggle(card)
  shift+click:   selection.range(focus.id, card)
  drag(card):    board.drag_card            # the gesture in `futures`
  alt+drag:      board.drag_copy
  shift+drag:    board.drag_card(axis: lock)
  wheel(board):  board.scroll(by: delta)
  rightclick:    menu.open(at: pointer, for: card)
  ctrl+z:        undo                       # runtime, not a block
\end{lstlisting}
\caption{\textbf{Desktop input is state plus bindings.} The input model of the desktop board of
\cref{sec:demos}: every key and mouse event names an action of some block, with parameters drawn from the
pointer, the focused item or the selection. Blocks never see a key.}
\label{fig:bindings}
\end{figure}

\subsection{Selection, focus, menus and undo}
\label{sec:selection}

Selection over $n$ items has $2^n$ states and would double the board's spec for every item; it is a block of its
own, read-wired to the board only through the item ids it names, and the board's actions take an item set as a
parameter. Focus is a block with one field. A context menu is a mode field with \texttt{open} and
\texttt{close} and a guard on every other action while it is open. Undo is not a block at all: it is a runtime
wrapper that keeps a history of machine states and restores one, which is exact by construction and costs
nothing to learn; learning it would multiply every block's state space by the depth of the history. Marquee
selection is a new kind of drag whose futures are selection sets: the candidates are the sets of items whose
rendered positions fall inside the rectangle swept so far, derived from the render block as other futures are
derived from the oracle, and it is the one piece of the desktop model that none of the Dragology demos we
catalogued contains. Text
entry stays out until the spec language has a text type.

\subsection{Drags inside a running machine}
\label{sec:livedrag}

A drag is a search over the futures the spec allows from the current state (\cref{sec:drags}). Inside a
running machine the grabbed block is detached from the clock for the duration of the gesture: events addressed
to it queue, the rest of the machine keeps running, and the pointer frame is minimised as before. On release
the drop is an external action and the queued events are delivered, so the cascade resumes from the dropped
state; the composed oracle says where the machine should be $k$ ticks later, and \cref{sec:demos} measures
whether it gets there. The visible signature of the whole design is that every block's spec check runs every
tick: a block that answered a legal transition lights green, one that did not lights red and the machine stops
there, with the transition on screen. Code shows nothing when it is right; this runtime shows its checks.

\subsection{Continuous hand-offs}
\label{sec:handoff}

Continuous blocks with one continuous field are within $1\%$ of range on $99.6$--$100\%$ of transitions, which
is not exact; one with three coupled fields (linked sliders, \cref{sec:round2}) is within tolerance on a tenth
of them; and the error compounds across steps: on the continuous slider a value that should stay at $0$ drifted to $0.59$ and then
$1.41$ over two clicks that the spec says do nothing (observed in the lab with the round-1 model;
\texttt{docs/paper2/timings.md}). A chain of continuous blocks would wander. The design rule
is that a hand-off is a discrete event: a continuous block emits an event only when a quantised condition holds
(the ball reached the end of the track), the event's payload is the spec's canonical value, and the receiving
block starts from that value, not from the sender's float. Inside a block the motion is continuous and may
carry its error for the length of one episode; across a hand-off the machine is exact again. The rolling-ball
experiment of \cref{sec:demos} measures drift with and without the rule.

\subsection{Rendering}
\label{sec:render}

The lab draws every element as a point, which was enough for a dial and a knob and is not enough for a
machine anyone would watch. The render block grows the shapes Part I's roadmap lists: rectangles, lines,
angles and z-order, each an expression over state, derived and never learned. The lab is the runtime: one page,
one tick loop, every block a network served from one process, with the Worker export of Part I as the way to
ship a machine once it runs. The networks are the behaviour; the picture is a function of the state, and we keep
it that way so that a machine can be checked without being looked at.
